\documentclass[11pt]{article}

    \usepackage[breakable]{tcolorbox}
    \usepackage{parskip} % Stop auto-indenting (to mimic markdown behaviour)
    

    % Basic figure setup, for now with no caption control since it's done
    % automatically by Pandoc (which extracts ![](path) syntax from Markdown).
    \usepackage{graphicx}
    % Keep aspect ratio if custom image width or height is specified
    \setkeys{Gin}{keepaspectratio}
    % Maintain compatibility with old templates. Remove in nbconvert 6.0
    \let\Oldincludegraphics\includegraphics
    % Ensure that by default, figures have no caption (until we provide a
    % proper Figure object with a Caption API and a way to capture that
    % in the conversion process - todo).
    \usepackage{caption}
    \DeclareCaptionFormat{nocaption}{}
    \captionsetup{format=nocaption,aboveskip=0pt,belowskip=0pt}

    \usepackage{float}
    \floatplacement{figure}{H} % forces figures to be placed at the correct location
    \usepackage{xcolor} % Allow colors to be defined
    \usepackage{enumerate} % Needed for markdown enumerations to work
    \usepackage{geometry} % Used to adjust the document margins
    \usepackage{amsmath} % Equations
    \usepackage{amssymb} % Equations
    \usepackage{textcomp} % defines textquotesingle
    % Hack from http://tex.stackexchange.com/a/47451/13684:
    \AtBeginDocument{%
        \def\PYZsq{\textquotesingle}% Upright quotes in Pygmentized code
    }
    \usepackage{upquote} % Upright quotes for verbatim code
    \usepackage{eurosym} % defines \euro

    \usepackage{iftex}
    \ifPDFTeX
        \usepackage[T1]{fontenc}
        \IfFileExists{alphabeta.sty}{
              \usepackage{alphabeta}
          }{
              \usepackage[mathletters]{ucs}
              \usepackage[utf8x]{inputenc}
          }
    \else
        \usepackage{fontspec}
        \usepackage{unicode-math}
    \fi

    \usepackage{fancyvrb} % verbatim replacement that allows latex
    \usepackage{grffile} % extends the file name processing of package graphics
                         % to support a larger range
    \makeatletter % fix for old versions of grffile with XeLaTeX
    \@ifpackagelater{grffile}{2019/11/01}
    {
      % Do nothing on new versions
    }
    {
      \def\Gread@@xetex#1{%
        \IfFileExists{"\Gin@base".bb}%
        {\Gread@eps{\Gin@base.bb}}%
        {\Gread@@xetex@aux#1}%
      }
    }
    \makeatother
    \usepackage[Export]{adjustbox} % Used to constrain images to a maximum size
    \adjustboxset{max size={0.9\linewidth}{0.9\paperheight}}

    % The hyperref package gives us a pdf with properly built
    % internal navigation ('pdf bookmarks' for the table of contents,
    % internal cross-reference links, web links for URLs, etc.)
    \usepackage{hyperref}
    % The default LaTeX title has an obnoxious amount of whitespace. By default,
    % titling removes some of it. It also provides customization options.
    \usepackage{titling}
    \usepackage{longtable} % longtable support required by pandoc >1.10
    \usepackage{booktabs}  % table support for pandoc > 1.12.2
    \usepackage{array}     % table support for pandoc >= 2.11.3
    \usepackage{calc}      % table minipage width calculation for pandoc >= 2.11.1
    \usepackage[inline]{enumitem} % IRkernel/repr support (it uses the enumerate* environment)
    \usepackage[normalem]{ulem} % ulem is needed to support strikethroughs (\sout)
                                % normalem makes italics be italics, not underlines
    \usepackage{soul}      % strikethrough (\st) support for pandoc >= 3.0.0
    \usepackage{mathrsfs}
    

    
    % Colors for the hyperref package
    \definecolor{urlcolor}{rgb}{0,.145,.698}
    \definecolor{linkcolor}{rgb}{.71,0.21,0.01}
    \definecolor{citecolor}{rgb}{.12,.54,.11}

    % ANSI colors
    \definecolor{ansi-black}{HTML}{3E424D}
    \definecolor{ansi-black-intense}{HTML}{282C36}
    \definecolor{ansi-red}{HTML}{E75C58}
    \definecolor{ansi-red-intense}{HTML}{B22B31}
    \definecolor{ansi-green}{HTML}{00A250}
    \definecolor{ansi-green-intense}{HTML}{007427}
    \definecolor{ansi-yellow}{HTML}{DDB62B}
    \definecolor{ansi-yellow-intense}{HTML}{B27D12}
    \definecolor{ansi-blue}{HTML}{208FFB}
    \definecolor{ansi-blue-intense}{HTML}{0065CA}
    \definecolor{ansi-magenta}{HTML}{D160C4}
    \definecolor{ansi-magenta-intense}{HTML}{A03196}
    \definecolor{ansi-cyan}{HTML}{60C6C8}
    \definecolor{ansi-cyan-intense}{HTML}{258F8F}
    \definecolor{ansi-white}{HTML}{C5C1B4}
    \definecolor{ansi-white-intense}{HTML}{A1A6B2}
    \definecolor{ansi-default-inverse-fg}{HTML}{FFFFFF}
    \definecolor{ansi-default-inverse-bg}{HTML}{000000}

    % common color for the border for error outputs.
    \definecolor{outerrorbackground}{HTML}{FFDFDF}

    % commands and environments needed by pandoc snippets
    % extracted from the output of `pandoc -s`
    \providecommand{\tightlist}{%
      \setlength{\itemsep}{0pt}\setlength{\parskip}{0pt}}
    \DefineVerbatimEnvironment{Highlighting}{Verbatim}{commandchars=\\\{\}}
    % Add ',fontsize=\small' for more characters per line
    \newenvironment{Shaded}{}{}
    \newcommand{\KeywordTok}[1]{\textcolor[rgb]{0.00,0.44,0.13}{\textbf{{#1}}}}
    \newcommand{\DataTypeTok}[1]{\textcolor[rgb]{0.56,0.13,0.00}{{#1}}}
    \newcommand{\DecValTok}[1]{\textcolor[rgb]{0.25,0.63,0.44}{{#1}}}
    \newcommand{\BaseNTok}[1]{\textcolor[rgb]{0.25,0.63,0.44}{{#1}}}
    \newcommand{\FloatTok}[1]{\textcolor[rgb]{0.25,0.63,0.44}{{#1}}}
    \newcommand{\CharTok}[1]{\textcolor[rgb]{0.25,0.44,0.63}{{#1}}}
    \newcommand{\StringTok}[1]{\textcolor[rgb]{0.25,0.44,0.63}{{#1}}}
    \newcommand{\CommentTok}[1]{\textcolor[rgb]{0.38,0.63,0.69}{\textit{{#1}}}}
    \newcommand{\OtherTok}[1]{\textcolor[rgb]{0.00,0.44,0.13}{{#1}}}
    \newcommand{\AlertTok}[1]{\textcolor[rgb]{1.00,0.00,0.00}{\textbf{{#1}}}}
    \newcommand{\FunctionTok}[1]{\textcolor[rgb]{0.02,0.16,0.49}{{#1}}}
    \newcommand{\RegionMarkerTok}[1]{{#1}}
    \newcommand{\ErrorTok}[1]{\textcolor[rgb]{1.00,0.00,0.00}{\textbf{{#1}}}}
    \newcommand{\NormalTok}[1]{{#1}}

    % Additional commands for more recent versions of Pandoc
    \newcommand{\ConstantTok}[1]{\textcolor[rgb]{0.53,0.00,0.00}{{#1}}}
    \newcommand{\SpecialCharTok}[1]{\textcolor[rgb]{0.25,0.44,0.63}{{#1}}}
    \newcommand{\VerbatimStringTok}[1]{\textcolor[rgb]{0.25,0.44,0.63}{{#1}}}
    \newcommand{\SpecialStringTok}[1]{\textcolor[rgb]{0.73,0.40,0.53}{{#1}}}
    \newcommand{\ImportTok}[1]{{#1}}
    \newcommand{\DocumentationTok}[1]{\textcolor[rgb]{0.73,0.13,0.13}{\textit{{#1}}}}
    \newcommand{\AnnotationTok}[1]{\textcolor[rgb]{0.38,0.63,0.69}{\textbf{\textit{{#1}}}}}
    \newcommand{\CommentVarTok}[1]{\textcolor[rgb]{0.38,0.63,0.69}{\textbf{\textit{{#1}}}}}
    \newcommand{\VariableTok}[1]{\textcolor[rgb]{0.10,0.09,0.49}{{#1}}}
    \newcommand{\ControlFlowTok}[1]{\textcolor[rgb]{0.00,0.44,0.13}{\textbf{{#1}}}}
    \newcommand{\OperatorTok}[1]{\textcolor[rgb]{0.40,0.40,0.40}{{#1}}}
    \newcommand{\BuiltInTok}[1]{{#1}}
    \newcommand{\ExtensionTok}[1]{{#1}}
    \newcommand{\PreprocessorTok}[1]{\textcolor[rgb]{0.74,0.48,0.00}{{#1}}}
    \newcommand{\AttributeTok}[1]{\textcolor[rgb]{0.49,0.56,0.16}{{#1}}}
    \newcommand{\InformationTok}[1]{\textcolor[rgb]{0.38,0.63,0.69}{\textbf{\textit{{#1}}}}}
    \newcommand{\WarningTok}[1]{\textcolor[rgb]{0.38,0.63,0.69}{\textbf{\textit{{#1}}}}}
    \makeatletter
    \newsavebox\pandoc@box
    \newcommand*\pandocbounded[1]{%
      \sbox\pandoc@box{#1}%
      % scaling factors for width and height
      \Gscale@div\@tempa\textheight{\dimexpr\ht\pandoc@box+\dp\pandoc@box\relax}%
      \Gscale@div\@tempb\linewidth{\wd\pandoc@box}%
      % select the smaller of both
      \ifdim\@tempb\p@<\@tempa\p@
        \let\@tempa\@tempb
      \fi
      % scaling accordingly (\@tempa < 1)
      \ifdim\@tempa\p@<\p@
        \scalebox{\@tempa}{\usebox\pandoc@box}%
      % scaling not needed, use as it is
      \else
        \usebox{\pandoc@box}%
      \fi
    }
    \makeatother

    % Define a nice break command that doesn't care if a line doesn't already
    % exist.
    \def\br{\hspace*{\fill} \\* }
    % Math Jax compatibility definitions
    \def\gt{>}
    \def\lt{<}
    \let\Oldtex\TeX
    \let\Oldlatex\LaTeX
    \renewcommand{\TeX}{\textrm{\Oldtex}}
    \renewcommand{\LaTeX}{\textrm{\Oldlatex}}
    % Document parameters
    % Document title
    \title{report}
    
    
    
    
    
    
    
% Pygments definitions
\makeatletter
\def\PY@reset{\let\PY@it=\relax \let\PY@bf=\relax%
    \let\PY@ul=\relax \let\PY@tc=\relax%
    \let\PY@bc=\relax \let\PY@ff=\relax}
\def\PY@tok#1{\csname PY@tok@#1\endcsname}
\def\PY@toks#1+{\ifx\relax#1\empty\else%
    \PY@tok{#1}\expandafter\PY@toks\fi}
\def\PY@do#1{\PY@bc{\PY@tc{\PY@ul{%
    \PY@it{\PY@bf{\PY@ff{#1}}}}}}}
\def\PY#1#2{\PY@reset\PY@toks#1+\relax+\PY@do{#2}}

\@namedef{PY@tok@w}{\def\PY@tc##1{\textcolor[rgb]{0.73,0.73,0.73}{##1}}}
\@namedef{PY@tok@c}{\let\PY@it=\textit\def\PY@tc##1{\textcolor[rgb]{0.24,0.48,0.48}{##1}}}
\@namedef{PY@tok@cp}{\def\PY@tc##1{\textcolor[rgb]{0.61,0.40,0.00}{##1}}}
\@namedef{PY@tok@k}{\let\PY@bf=\textbf\def\PY@tc##1{\textcolor[rgb]{0.00,0.50,0.00}{##1}}}
\@namedef{PY@tok@kp}{\def\PY@tc##1{\textcolor[rgb]{0.00,0.50,0.00}{##1}}}
\@namedef{PY@tok@kt}{\def\PY@tc##1{\textcolor[rgb]{0.69,0.00,0.25}{##1}}}
\@namedef{PY@tok@o}{\def\PY@tc##1{\textcolor[rgb]{0.40,0.40,0.40}{##1}}}
\@namedef{PY@tok@ow}{\let\PY@bf=\textbf\def\PY@tc##1{\textcolor[rgb]{0.67,0.13,1.00}{##1}}}
\@namedef{PY@tok@nb}{\def\PY@tc##1{\textcolor[rgb]{0.00,0.50,0.00}{##1}}}
\@namedef{PY@tok@nf}{\def\PY@tc##1{\textcolor[rgb]{0.00,0.00,1.00}{##1}}}
\@namedef{PY@tok@nc}{\let\PY@bf=\textbf\def\PY@tc##1{\textcolor[rgb]{0.00,0.00,1.00}{##1}}}
\@namedef{PY@tok@nn}{\let\PY@bf=\textbf\def\PY@tc##1{\textcolor[rgb]{0.00,0.00,1.00}{##1}}}
\@namedef{PY@tok@ne}{\let\PY@bf=\textbf\def\PY@tc##1{\textcolor[rgb]{0.80,0.25,0.22}{##1}}}
\@namedef{PY@tok@nv}{\def\PY@tc##1{\textcolor[rgb]{0.10,0.09,0.49}{##1}}}
\@namedef{PY@tok@no}{\def\PY@tc##1{\textcolor[rgb]{0.53,0.00,0.00}{##1}}}
\@namedef{PY@tok@nl}{\def\PY@tc##1{\textcolor[rgb]{0.46,0.46,0.00}{##1}}}
\@namedef{PY@tok@ni}{\let\PY@bf=\textbf\def\PY@tc##1{\textcolor[rgb]{0.44,0.44,0.44}{##1}}}
\@namedef{PY@tok@na}{\def\PY@tc##1{\textcolor[rgb]{0.41,0.47,0.13}{##1}}}
\@namedef{PY@tok@nt}{\let\PY@bf=\textbf\def\PY@tc##1{\textcolor[rgb]{0.00,0.50,0.00}{##1}}}
\@namedef{PY@tok@nd}{\def\PY@tc##1{\textcolor[rgb]{0.67,0.13,1.00}{##1}}}
\@namedef{PY@tok@s}{\def\PY@tc##1{\textcolor[rgb]{0.73,0.13,0.13}{##1}}}
\@namedef{PY@tok@sd}{\let\PY@it=\textit\def\PY@tc##1{\textcolor[rgb]{0.73,0.13,0.13}{##1}}}
\@namedef{PY@tok@si}{\let\PY@bf=\textbf\def\PY@tc##1{\textcolor[rgb]{0.64,0.35,0.47}{##1}}}
\@namedef{PY@tok@se}{\let\PY@bf=\textbf\def\PY@tc##1{\textcolor[rgb]{0.67,0.36,0.12}{##1}}}
\@namedef{PY@tok@sr}{\def\PY@tc##1{\textcolor[rgb]{0.64,0.35,0.47}{##1}}}
\@namedef{PY@tok@ss}{\def\PY@tc##1{\textcolor[rgb]{0.10,0.09,0.49}{##1}}}
\@namedef{PY@tok@sx}{\def\PY@tc##1{\textcolor[rgb]{0.00,0.50,0.00}{##1}}}
\@namedef{PY@tok@m}{\def\PY@tc##1{\textcolor[rgb]{0.40,0.40,0.40}{##1}}}
\@namedef{PY@tok@gh}{\let\PY@bf=\textbf\def\PY@tc##1{\textcolor[rgb]{0.00,0.00,0.50}{##1}}}
\@namedef{PY@tok@gu}{\let\PY@bf=\textbf\def\PY@tc##1{\textcolor[rgb]{0.50,0.00,0.50}{##1}}}
\@namedef{PY@tok@gd}{\def\PY@tc##1{\textcolor[rgb]{0.63,0.00,0.00}{##1}}}
\@namedef{PY@tok@gi}{\def\PY@tc##1{\textcolor[rgb]{0.00,0.52,0.00}{##1}}}
\@namedef{PY@tok@gr}{\def\PY@tc##1{\textcolor[rgb]{0.89,0.00,0.00}{##1}}}
\@namedef{PY@tok@ge}{\let\PY@it=\textit}
\@namedef{PY@tok@gs}{\let\PY@bf=\textbf}
\@namedef{PY@tok@ges}{\let\PY@bf=\textbf\let\PY@it=\textit}
\@namedef{PY@tok@gp}{\let\PY@bf=\textbf\def\PY@tc##1{\textcolor[rgb]{0.00,0.00,0.50}{##1}}}
\@namedef{PY@tok@go}{\def\PY@tc##1{\textcolor[rgb]{0.44,0.44,0.44}{##1}}}
\@namedef{PY@tok@gt}{\def\PY@tc##1{\textcolor[rgb]{0.00,0.27,0.87}{##1}}}
\@namedef{PY@tok@err}{\def\PY@bc##1{{\setlength{\fboxsep}{\string -\fboxrule}\fcolorbox[rgb]{1.00,0.00,0.00}{1,1,1}{\strut ##1}}}}
\@namedef{PY@tok@kc}{\let\PY@bf=\textbf\def\PY@tc##1{\textcolor[rgb]{0.00,0.50,0.00}{##1}}}
\@namedef{PY@tok@kd}{\let\PY@bf=\textbf\def\PY@tc##1{\textcolor[rgb]{0.00,0.50,0.00}{##1}}}
\@namedef{PY@tok@kn}{\let\PY@bf=\textbf\def\PY@tc##1{\textcolor[rgb]{0.00,0.50,0.00}{##1}}}
\@namedef{PY@tok@kr}{\let\PY@bf=\textbf\def\PY@tc##1{\textcolor[rgb]{0.00,0.50,0.00}{##1}}}
\@namedef{PY@tok@bp}{\def\PY@tc##1{\textcolor[rgb]{0.00,0.50,0.00}{##1}}}
\@namedef{PY@tok@fm}{\def\PY@tc##1{\textcolor[rgb]{0.00,0.00,1.00}{##1}}}
\@namedef{PY@tok@vc}{\def\PY@tc##1{\textcolor[rgb]{0.10,0.09,0.49}{##1}}}
\@namedef{PY@tok@vg}{\def\PY@tc##1{\textcolor[rgb]{0.10,0.09,0.49}{##1}}}
\@namedef{PY@tok@vi}{\def\PY@tc##1{\textcolor[rgb]{0.10,0.09,0.49}{##1}}}
\@namedef{PY@tok@vm}{\def\PY@tc##1{\textcolor[rgb]{0.10,0.09,0.49}{##1}}}
\@namedef{PY@tok@sa}{\def\PY@tc##1{\textcolor[rgb]{0.73,0.13,0.13}{##1}}}
\@namedef{PY@tok@sb}{\def\PY@tc##1{\textcolor[rgb]{0.73,0.13,0.13}{##1}}}
\@namedef{PY@tok@sc}{\def\PY@tc##1{\textcolor[rgb]{0.73,0.13,0.13}{##1}}}
\@namedef{PY@tok@dl}{\def\PY@tc##1{\textcolor[rgb]{0.73,0.13,0.13}{##1}}}
\@namedef{PY@tok@s2}{\def\PY@tc##1{\textcolor[rgb]{0.73,0.13,0.13}{##1}}}
\@namedef{PY@tok@sh}{\def\PY@tc##1{\textcolor[rgb]{0.73,0.13,0.13}{##1}}}
\@namedef{PY@tok@s1}{\def\PY@tc##1{\textcolor[rgb]{0.73,0.13,0.13}{##1}}}
\@namedef{PY@tok@mb}{\def\PY@tc##1{\textcolor[rgb]{0.40,0.40,0.40}{##1}}}
\@namedef{PY@tok@mf}{\def\PY@tc##1{\textcolor[rgb]{0.40,0.40,0.40}{##1}}}
\@namedef{PY@tok@mh}{\def\PY@tc##1{\textcolor[rgb]{0.40,0.40,0.40}{##1}}}
\@namedef{PY@tok@mi}{\def\PY@tc##1{\textcolor[rgb]{0.40,0.40,0.40}{##1}}}
\@namedef{PY@tok@il}{\def\PY@tc##1{\textcolor[rgb]{0.40,0.40,0.40}{##1}}}
\@namedef{PY@tok@mo}{\def\PY@tc##1{\textcolor[rgb]{0.40,0.40,0.40}{##1}}}
\@namedef{PY@tok@ch}{\let\PY@it=\textit\def\PY@tc##1{\textcolor[rgb]{0.24,0.48,0.48}{##1}}}
\@namedef{PY@tok@cm}{\let\PY@it=\textit\def\PY@tc##1{\textcolor[rgb]{0.24,0.48,0.48}{##1}}}
\@namedef{PY@tok@cpf}{\let\PY@it=\textit\def\PY@tc##1{\textcolor[rgb]{0.24,0.48,0.48}{##1}}}
\@namedef{PY@tok@c1}{\let\PY@it=\textit\def\PY@tc##1{\textcolor[rgb]{0.24,0.48,0.48}{##1}}}
\@namedef{PY@tok@cs}{\let\PY@it=\textit\def\PY@tc##1{\textcolor[rgb]{0.24,0.48,0.48}{##1}}}

\def\PYZbs{\char`\\}
\def\PYZus{\char`\_}
\def\PYZob{\char`\{}
\def\PYZcb{\char`\}}
\def\PYZca{\char`\^}
\def\PYZam{\char`\&}
\def\PYZlt{\char`\<}
\def\PYZgt{\char`\>}
\def\PYZsh{\char`\#}
\def\PYZpc{\char`\%}
\def\PYZdl{\char`\$}
\def\PYZhy{\char`\-}
\def\PYZsq{\char`\'}
\def\PYZdq{\char`\"}
\def\PYZti{\char`\~}
% for compatibility with earlier versions
\def\PYZat{@}
\def\PYZlb{[}
\def\PYZrb{]}
\makeatother


    % For linebreaks inside Verbatim environment from package fancyvrb.
    \makeatletter
        \newbox\Wrappedcontinuationbox
        \newbox\Wrappedvisiblespacebox
        \newcommand*\Wrappedvisiblespace {\textcolor{red}{\textvisiblespace}}
        \newcommand*\Wrappedcontinuationsymbol {\textcolor{red}{\llap{\tiny$\m@th\hookrightarrow$}}}
        \newcommand*\Wrappedcontinuationindent {3ex }
        \newcommand*\Wrappedafterbreak {\kern\Wrappedcontinuationindent\copy\Wrappedcontinuationbox}
        % Take advantage of the already applied Pygments mark-up to insert
        % potential linebreaks for TeX processing.
        %        {, <, #, %, $, ' and ": go to next line.
        %        _, }, ^, &, >, - and ~: stay at end of broken line.
        % Use of \textquotesingle for straight quote.
        \newcommand*\Wrappedbreaksatspecials {%
            \def\PYGZus{\discretionary{\char`\_}{\Wrappedafterbreak}{\char`\_}}%
            \def\PYGZob{\discretionary{}{\Wrappedafterbreak\char`\{}{\char`\{}}%
            \def\PYGZcb{\discretionary{\char`\}}{\Wrappedafterbreak}{\char`\}}}%
            \def\PYGZca{\discretionary{\char`\^}{\Wrappedafterbreak}{\char`\^}}%
            \def\PYGZam{\discretionary{\char`\&}{\Wrappedafterbreak}{\char`\&}}%
            \def\PYGZlt{\discretionary{}{\Wrappedafterbreak\char`\<}{\char`\<}}%
            \def\PYGZgt{\discretionary{\char`\>}{\Wrappedafterbreak}{\char`\>}}%
            \def\PYGZsh{\discretionary{}{\Wrappedafterbreak\char`\#}{\char`\#}}%
            \def\PYGZpc{\discretionary{}{\Wrappedafterbreak\char`\%}{\char`\%}}%
            \def\PYGZdl{\discretionary{}{\Wrappedafterbreak\char`\$}{\char`\$}}%
            \def\PYGZhy{\discretionary{\char`\-}{\Wrappedafterbreak}{\char`\-}}%
            \def\PYGZsq{\discretionary{}{\Wrappedafterbreak\textquotesingle}{\textquotesingle}}%
            \def\PYGZdq{\discretionary{}{\Wrappedafterbreak\char`\"}{\char`\"}}%
            \def\PYGZti{\discretionary{\char`\~}{\Wrappedafterbreak}{\char`\~}}%
        }
        % Some characters . , ; ? ! / are not pygmentized.
        % This macro makes them "active" and they will insert potential linebreaks
        \newcommand*\Wrappedbreaksatpunct {%
            \lccode`\~`\.\lowercase{\def~}{\discretionary{\hbox{\char`\.}}{\Wrappedafterbreak}{\hbox{\char`\.}}}%
            \lccode`\~`\,\lowercase{\def~}{\discretionary{\hbox{\char`\,}}{\Wrappedafterbreak}{\hbox{\char`\,}}}%
            \lccode`\~`\;\lowercase{\def~}{\discretionary{\hbox{\char`\;}}{\Wrappedafterbreak}{\hbox{\char`\;}}}%
            \lccode`\~`\:\lowercase{\def~}{\discretionary{\hbox{\char`\:}}{\Wrappedafterbreak}{\hbox{\char`\:}}}%
            \lccode`\~`\?\lowercase{\def~}{\discretionary{\hbox{\char`\?}}{\Wrappedafterbreak}{\hbox{\char`\?}}}%
            \lccode`\~`\!\lowercase{\def~}{\discretionary{\hbox{\char`\!}}{\Wrappedafterbreak}{\hbox{\char`\!}}}%
            \lccode`\~`\/\lowercase{\def~}{\discretionary{\hbox{\char`\/}}{\Wrappedafterbreak}{\hbox{\char`\/}}}%
            \catcode`\.\active
            \catcode`\,\active
            \catcode`\;\active
            \catcode`\:\active
            \catcode`\?\active
            \catcode`\!\active
            \catcode`\/\active
            \lccode`\~`\~
        }
    \makeatother

    \let\OriginalVerbatim=\Verbatim
    \makeatletter
    \renewcommand{\Verbatim}[1][1]{%
        %\parskip\z@skip
        \sbox\Wrappedcontinuationbox {\Wrappedcontinuationsymbol}%
        \sbox\Wrappedvisiblespacebox {\FV@SetupFont\Wrappedvisiblespace}%
        \def\FancyVerbFormatLine ##1{\hsize\linewidth
            \vtop{\raggedright\hyphenpenalty\z@\exhyphenpenalty\z@
                \doublehyphendemerits\z@\finalhyphendemerits\z@
                \strut ##1\strut}%
        }%
        % If the linebreak is at a space, the latter will be displayed as visible
        % space at end of first line, and a continuation symbol starts next line.
        % Stretch/shrink are however usually zero for typewriter font.
        \def\FV@Space {%
            \nobreak\hskip\z@ plus\fontdimen3\font minus\fontdimen4\font
            \discretionary{\copy\Wrappedvisiblespacebox}{\Wrappedafterbreak}
            {\kern\fontdimen2\font}%
        }%

        % Allow breaks at special characters using \PYG... macros.
        \Wrappedbreaksatspecials
        % Breaks at punctuation characters . , ; ? ! and / need catcode=\active
        \OriginalVerbatim[#1,codes*=\Wrappedbreaksatpunct]%
    }
    \makeatother

    % Exact colors from NB
    \definecolor{incolor}{HTML}{303F9F}
    \definecolor{outcolor}{HTML}{D84315}
    \definecolor{cellborder}{HTML}{CFCFCF}
    \definecolor{cellbackground}{HTML}{F7F7F7}

    % prompt
    \makeatletter
    \newcommand{\boxspacing}{\kern\kvtcb@left@rule\kern\kvtcb@boxsep}
    \makeatother
    \newcommand{\prompt}[4]{
        {\ttfamily\llap{{\color{#2}[#3]:\hspace{3pt}#4}}\vspace{-\baselineskip}}
    }
    

    
    % Prevent overflowing lines due to hard-to-break entities
    \sloppy
    % Setup hyperref package
    \hypersetup{
      breaklinks=true,  % so long urls are correctly broken across lines
      colorlinks=true,
      urlcolor=urlcolor,
      linkcolor=linkcolor,
      citecolor=citecolor,
      }
    % Slightly bigger margins than the latex defaults
    
    \geometry{verbose,tmargin=1in,bmargin=1in,lmargin=1in,rmargin=1in}
    
    

\begin{document}
    
    \maketitle
    
    

    
    \subsection{Overview}\label{overview}

These notes summarize the replication of \textbf{Figure 1} and
\textbf{Tables 1, 3, and 4}.

    \subsection{Data}\label{data}

Monthly with and without dividend returns of the value weighted
portfolio of all NYSE, Amex, and NASDAQ stocks from 1944 to 2007 and
2024.

30-day T-bill rate from 1944 to 2007 and 2024.

The reason for needing the data from 1944 instead of 1946 was due to the
construction of certain parameters requiring data from previous years.

    \subsection{Notation}\label{notation}

For month \(m\),

\begin{itemize}
\tightlist
\item
  \(r_m^d\) = monthly value-weighted return \textbf{with} dividends
  (\texttt{vwretd}),
\item
  \(r_m^x\) = monthly value-weighted return \textbf{without} dividends
  (\texttt{vwretx}),
\item
  \(r_m^f\) = monthly T-bill return (\texttt{t30ret}),
\item
  \(P_m\) = ex-dividend price index at month \(m\),
\item
  \(\mathrm{Div}_m\) = dividend cash flow paid during month \(m\).
\end{itemize}

    \subsection{Monthly-to-Annual
Construction}\label{monthly-to-annual-construction}

\subsubsection{1. Ex-dividend price
index}\label{ex-dividend-price-index}

The ex-dividend return satisfies

\[
1+r_m^x = \frac{P_m}{P_{m-1}}.
\]

So the ex-dividend price index is built recursively as

\[
P_m = P_{m-1}(1+r_m^x),
\]

starting from an arbitrary normalization such as \(P_0=1\).

\subsubsection{2. Monthly dividend cash
flow}\label{monthly-dividend-cash-flow}

The cum-dividend return satisfies

\[
1+r_m^d = \frac{P_m+\mathrm{Div}_m}{P_{m-1}}.
\]

Subtracting the ex-dividend equation gives

\[
\mathrm{Div}_m = P_{m-1}\bigl(r_m^d-r_m^x\bigr).
\]

This is the monthly dividend cash flow.

\subsubsection{3. Annual dividends under two reinvestment
conventions}\label{annual-dividends-under-two-reinvestment-conventions}

For each calendar year \(t\), the monthly dividends are carried to the
end of that year in two different ways.

\paragraph{Risk-free reinvestment}\label{risk-free-reinvestment}

\[
D_t
=
\sum_{m \in t}
\mathrm{Div}_m
\prod_{\substack{j>m \\ j \in t}}
(1+r_j^f).
\]

\paragraph{Market reinvestment}\label{market-reinvestment}

\[
D_t^M
=
\sum_{m \in t}
\mathrm{Div}_m
\prod_{\substack{j>m \\ j \in t}}
(1+r_j^d).
\]

\subsubsection{4. Annual dividend growth}\label{annual-dividend-growth}

The annual log dividend-growth series are

\[
\Delta d_t
=
\log\left(\frac{D_t}{D_{t-1}}\right),
\qquad
\Delta d_t^M
=
\log\left(\frac{D_t^M}{D_{t-1}^M}\right).
\]

These two series are the objects behind both \textbf{Figure 1} and
\textbf{Table 1}.

\subsubsection{5. Annual returns and log price-dividend
ratios}\label{annual-returns-and-log-price-dividend-ratios}

Let \(P_t\) denote the end-of-year ex-dividend price index. Then

\[
r_t
=
\log\left(\frac{P_t+D_t}{P_{t-1}}\right),
\qquad
r_t^M
=
\log\left(\frac{P_t+D_t^M}{P_{t-1}}\right).
\]

The annual log price-dividend ratios are

\[
pd_t = \log\left(\frac{P_t}{D_t}\right),
\qquad
pd_t^M = \log\left(\frac{P_t}{D_t^M}\right).
\]

    \subsection{Figure 1}\label{figure-1}

Figure 1 is a time-series plot of

\[
\Delta d_t
\quad \text{and} \quad
\Delta d_t^M
\]

over the annual sample.

Replication was successful.

\subsubsection{Replicated Figure 1}\label{replicated-figure-1}

\begin{figure}
\centering
\pandocbounded{\includegraphics[keepaspectratio,alt={2007}]{/Users/moxiaoli/Desktop/Uchicago/2025 - 2026/p03_binsbergen_koijen_2010/_output/figure1_2007.png}}
\caption{2007}
\end{figure}

\subsubsection{Paper Figure 1}\label{paper-figure-1}

\begin{figure}
\centering
\pandocbounded{\includegraphics[keepaspectratio,alt={paper}]{/Users/moxiaoli/Desktop/Uchicago/2025 - 2026/p03_binsbergen_koijen_2010/_output/figure1.png}}
\caption{paper}
\end{figure}

\subsubsection{Replicated Figure 1 (1946 -
2024)}\label{replicated-figure-1-1946---2024}

\begin{figure}
\centering
\pandocbounded{\includegraphics[keepaspectratio,alt={2007}]{/Users/moxiaoli/Desktop/Uchicago/2025 - 2026/p03_binsbergen_koijen_2010/_output/figure1_2024.png}}
\caption{2007}
\end{figure}

    \subsection{Table 1}\label{table-1}

Table 1 reports summary statistics of the two annual dividend-growth
series.

Replication was successful.

\subsubsection{Replicated Table 1}\label{replicated-table-1}

{\def\LTcaptype{none} % do not increment counter
\begin{longtable}[]{@{}lrr@{}}
\toprule\noalign{}
Statistic & \(\Delta d_t^M\) & \(\Delta d_t\) \\
\midrule\noalign{}
\endhead
\bottomrule\noalign{}
\endlastfoot
Mean & 0.0588 & 0.0613 \\
Median & 0.0534 & 0.0544 \\
Standard deviation & 0.1237 & 0.0624 \\
Maximum & 0.3697 & 0.2609 \\
Minimum & -0.2884 & -0.0544 \\
\end{longtable}
}

\subsubsection{Paper Table 1}\label{paper-table-1}

{\def\LTcaptype{none} % do not increment counter
\begin{longtable}[]{@{}lrr@{}}
\toprule\noalign{}
Statistic & \(\Delta d_t^M\) & \(\Delta d_t\) \\
\midrule\noalign{}
\endhead
\bottomrule\noalign{}
\endlastfoot
Mean & 0.0586 & 0.0611 \\
Median & 0.0558 & 0.0540 \\
Standard deviation & 0.1232 & 0.0622 \\
Maximum & 0.3699 & 0.2616 \\
Minimum & -0.2912 & -0.0579 \\
\end{longtable}
}

\subsubsection{Side-by-side comparison}\label{side-by-side-comparison}

{\def\LTcaptype{none} % do not increment counter
\begin{longtable}[]{@{}
  >{\raggedright\arraybackslash}p{(\linewidth - 8\tabcolsep) * \real{0.1579}}
  >{\raggedleft\arraybackslash}p{(\linewidth - 8\tabcolsep) * \real{0.2105}}
  >{\raggedleft\arraybackslash}p{(\linewidth - 8\tabcolsep) * \real{0.2105}}
  >{\raggedleft\arraybackslash}p{(\linewidth - 8\tabcolsep) * \real{0.2105}}
  >{\raggedleft\arraybackslash}p{(\linewidth - 8\tabcolsep) * \real{0.2105}}@{}}
\toprule\noalign{}
\begin{minipage}[b]{\linewidth}\raggedright
Statistic
\end{minipage} & \begin{minipage}[b]{\linewidth}\raggedleft
Replicated \(\Delta d_t^M\)
\end{minipage} & \begin{minipage}[b]{\linewidth}\raggedleft
Paper \(\Delta d_t^M\)
\end{minipage} & \begin{minipage}[b]{\linewidth}\raggedleft
Replicated \(\Delta d_t\)
\end{minipage} & \begin{minipage}[b]{\linewidth}\raggedleft
Paper \(\Delta d_t\)
\end{minipage} \\
\midrule\noalign{}
\endhead
\bottomrule\noalign{}
\endlastfoot
Mean & 0.0588 & 0.0586 & 0.0613 & 0.0611 \\
Median & 0.0534 & 0.0558 & 0.0544 & 0.0540 \\
Standard deviation & 0.1237 & 0.1232 & 0.0624 & 0.0622 \\
Maximum & 0.3697 & 0.3699 & 0.2609 & 0.2616 \\
Minimum & -0.2884 & -0.2912 & -0.0544 & -0.0579 \\
\end{longtable}
}

\subsubsection{Replicated Table 1 (1946 -
2024)}\label{replicated-table-1-1946---2024}

{\def\LTcaptype{none} % do not increment counter
\begin{longtable}[]{@{}lrr@{}}
\toprule\noalign{}
Statistic & \(\Delta d_t^M\) & \(\Delta d_t\) \\
\midrule\noalign{}
\endhead
\bottomrule\noalign{}
\endlastfoot
Mean & 0.0587 & 0.0599 \\
Median & 0.0534 & 0.0547 \\
Standard deviation & 0.1284 & 0.0678 \\
Maximum & 0.3697 & 0.2609 \\
Minimum & -0.2884 & -0.1885 \\
\end{longtable}
}

    \subsection{Table 3}\label{table-3}

Table 3 uses predictive OLS of annual returns and annual dividend growth
on lagged log price-dividend ratios.

The four regressions are

\[
r_t^M = \alpha_r^M + \beta_r^M pd_{t-1}^M + \varepsilon_t^{r,M},
\]

\[
r_t = \alpha_r + \beta_r pd_{t-1} + \varepsilon_t^{r},
\]

\[
\Delta d_t^M = \alpha_d^M + \beta_d^M pd_{t-1}^M + \varepsilon_t^{d,M},
\]

\[
\Delta d_t = \alpha_d + \beta_d pd_{t-1} + \varepsilon_t^{d}.
\]

The replication was successful.

\subsubsection{Replicated Table 3}\label{replicated-table-3}

{\def\LTcaptype{none} % do not increment counter
\begin{longtable}[]{@{}lrrrr@{}}
\toprule\noalign{}
& \(r_t^M\) & \(r_t\) & \(\Delta d_t^M\) & \(\Delta d_t\) \\
\midrule\noalign{}
\endhead
\bottomrule\noalign{}
\endlastfoot
Constant & 0.4572 & 0.4574 & 0.1812 & 0.1094 \\
SE(Constant) & 0.1541 & 0.1528 & 0.1273 & 0.0648 \\
\(pd_{t-1}\) or \(pd_{t-1}^M\) & -0.1032 & -0.1026 & -0.0360 &
-0.0140 \\
SE(\(pd_{t-1}\) or \(pd_{t-1}^M\)) & 0.0450 & 0.0442 & 0.0372 &
0.0187 \\
\(R^2\) & 0.0806 & 0.0825 & 0.0154 & 0.0092 \\
Adj. \(R^2\) & 0.0653 & 0.0672 & -0.0010 & -0.0073 \\
\end{longtable}
}

\subsubsection{Paper Table 3}\label{paper-table-3}

{\def\LTcaptype{none} % do not increment counter
\begin{longtable}[]{@{}lrrrr@{}}
\toprule\noalign{}
& \(r_t^M\) & \(r_t\) & \(\Delta d_t^M\) & \(\Delta d_t\) \\
\midrule\noalign{}
\endhead
\bottomrule\noalign{}
\endlastfoot
Constant & 0.4539 & 0.4555 & 0.1814 & 0.1085 \\
SE(Constant) & 0.1537 & 0.1524 & 0.1266 & 0.0645 \\
\(pd_{t-1}\) or \(pd_{t-1}^M\) & -0.1023 & -0.1020 & -0.0361 &
-0.0138 \\
SE(\(pd_{t-1}\) or \(pd_{t-1}^M\)) & 0.0449 & 0.0441 & 0.0370 &
0.0186 \\
\(R^2\) & 0.0796 & 0.0820 & 0.0156 & 0.0090 \\
Adj. \(R^2\) & 0.0643 & 0.0667 & -0.0007 & -0.0075 \\
\end{longtable}
}

\subsubsection{Side-by-side comparison}\label{side-by-side-comparison}

{\def\LTcaptype{none} % do not increment counter
\begin{longtable}[]{@{}
  >{\raggedright\arraybackslash}p{(\linewidth - 16\tabcolsep) * \real{0.0857}}
  >{\raggedleft\arraybackslash}p{(\linewidth - 16\tabcolsep) * \real{0.1143}}
  >{\raggedleft\arraybackslash}p{(\linewidth - 16\tabcolsep) * \real{0.1143}}
  >{\raggedleft\arraybackslash}p{(\linewidth - 16\tabcolsep) * \real{0.1143}}
  >{\raggedleft\arraybackslash}p{(\linewidth - 16\tabcolsep) * \real{0.1143}}
  >{\raggedleft\arraybackslash}p{(\linewidth - 16\tabcolsep) * \real{0.1143}}
  >{\raggedleft\arraybackslash}p{(\linewidth - 16\tabcolsep) * \real{0.1143}}
  >{\raggedleft\arraybackslash}p{(\linewidth - 16\tabcolsep) * \real{0.1143}}
  >{\raggedleft\arraybackslash}p{(\linewidth - 16\tabcolsep) * \real{0.1143}}@{}}
\toprule\noalign{}
\begin{minipage}[b]{\linewidth}\raggedright
\end{minipage} & \begin{minipage}[b]{\linewidth}\raggedleft
Rep.~\(r_t^M\)
\end{minipage} & \begin{minipage}[b]{\linewidth}\raggedleft
Paper \(r_t^M\)
\end{minipage} & \begin{minipage}[b]{\linewidth}\raggedleft
Rep.~\(r_t\)
\end{minipage} & \begin{minipage}[b]{\linewidth}\raggedleft
Paper \(r_t\)
\end{minipage} & \begin{minipage}[b]{\linewidth}\raggedleft
Rep.~\(\Delta d_t^M\)
\end{minipage} & \begin{minipage}[b]{\linewidth}\raggedleft
Paper \(\Delta d_t^M\)
\end{minipage} & \begin{minipage}[b]{\linewidth}\raggedleft
Rep.~\(\Delta d_t\)
\end{minipage} & \begin{minipage}[b]{\linewidth}\raggedleft
Paper \(\Delta d_t\)
\end{minipage} \\
\midrule\noalign{}
\endhead
\bottomrule\noalign{}
\endlastfoot
Constant & 0.4572 & 0.4539 & 0.4574 & 0.4555 & 0.1812 & 0.1814 & 0.1094
& 0.1085 \\
SE(Constant) & 0.1541 & 0.1537 & 0.1528 & 0.1524 & 0.1273 & 0.1266 &
0.0648 & 0.0645 \\
\(pd_{t-1}\) or \(pd_{t-1}^M\) & -0.1032 & -0.1023 & -0.1026 & -0.1020 &
-0.0360 & -0.0361 & -0.0140 & -0.0138 \\
SE(\(pd_{t-1}\) or \(pd_{t-1}^M\)) & 0.0450 & 0.0449 & 0.0442 & 0.0441 &
0.0372 & 0.0370 & 0.0187 & 0.0186 \\
\(R^2\) & 0.0806 & 0.0796 & 0.0825 & 0.0820 & 0.0154 & 0.0156 & 0.0092 &
0.0090 \\
Adj. \(R^2\) & 0.0653 & 0.0643 & 0.0672 & 0.0667 & -0.0010 & -0.0007 &
-0.0073 & -0.0075 \\
\end{longtable}
}

\subsubsection{Replicated Table 3 (1946 -
2024)}\label{replicated-table-3-1946---2024}

{\def\LTcaptype{none} % do not increment counter
\begin{longtable}[]{@{}lrrrr@{}}
\toprule\noalign{}
& \(r_t^M\) & \(r_t\) & \(\Delta d_t^M\) & \(\Delta d_t\) \\
\midrule\noalign{}
\endhead
\bottomrule\noalign{}
\endlastfoot
Constant & 0.4525 & 0.4580 & 0.1772 & 0.0765 \\
SE(Constant) & 0.1482 & 0.1453 & 0.1194 & 0.0632 \\
\(pd_{t-1}\) or \(pd_{t-1}^M\) & -0.0998 & -0.1007 & -0.0339 &
-0.0047 \\
SE(\(pd_{t-1}\) or \(pd_{t-1}^M\)) & 0.0420 & 0.0408 & 0.0338 &
0.0178 \\
\(R^2\) & 0.0684 & 0.0733 & 0.0128 & 0.0009 \\
Adj. \(R^2\) & 0.0563 & 0.0613 & 0.0000 & -0.0121 \\
\end{longtable}
}

    \subsection{Table 4}\label{table-4}

Table 4 studies reduced-form time-series models for market-reinvested
dividend growth. Let

\[
y_t := \Delta d_t^M.
\]

\subsubsection{Column 1: AR(1)}\label{column-1-ar1}

\[
y_t = a_1 + b_1 y_{t-1} + u_t.
\]

This was estimated by OLS.

\subsubsection{Column 2: lagged-return
regression}\label{column-2-lagged-return-regression}

\[
y_t = a_2 + b_2 r_{t-1} + u_t.
\]

This was also estimated by OLS.

\subsubsection{Column 3: Initial ARIMA
attempt}\label{column-3-initial-arima-attempt}

The initial attempt used the ARIMA package, which implements a generic
ARIMA model with an exogenous regressor:

\[
y_t
=
(1-\phi)c
+ \phi y_{t-1}
+ \beta x_t
- \phi \beta x_{t-1}
+ e_t + \theta e_{t-1}.
\]

Under this approach, the resulting column-3 replications were
unsatisfactory. \#\#\# Column 3: Direct ARMAX replacement

To fix this, the ARIMA-with-exogenous-regressor approach was replaced by
a \textbf{direct estimation}.

The final specification was estimated as

\[
y_t = c + \phi y_{t-1} + \beta r_{t-1} + e_t + \theta e_{t-1},
\]

with \(e_0=0\).

Rearranging gives the innovation recursion

\[
e_t = y_t - c - \phi y_{t-1} - \beta r_{t-1} - \theta e_{t-1}.
\]

For a candidate parameter vector

\[
\psi = (c,\phi,\theta,\beta),
\]

by MLE, the conditional innovation variance is

\[
\hat{\sigma}^2(\psi) = \frac{1}{T}\sum_{t=1}^T e_t(\psi)^2,
\]

and the conditional Gaussian negative log-likelihood minimized in the
replication is

\[
Q(\psi)
=
\frac{T}{2}
\left[
\log(2\pi)+1+\log \hat{\sigma}^2(\psi)
\right],
\]

which is minimized by SciPy.

\subsubsection{Reported constant in column
3}\label{reported-constant-in-column-3}

The raw intercept is \(c\), but the reported column-3 constant was taken
to be the unconditional mean

\[
\gamma_0
=
\frac{c+\beta \bar r}{1-\phi},
\qquad
\bar r = \frac{1}{T}\sum_{t=1}^T r_{t-1}.
\]

Its standard error was computed using the delta method with a covariance
matrix based on the numerical Hessian of the objective.

Replication was successful.

\subsubsection{Replicated Table 4}\label{replicated-table-4}

{\def\LTcaptype{none} % do not increment counter
\begin{longtable}[]{@{}lrrr@{}}
\toprule\noalign{}
& (1) AR(1) & (2) \(r_{t-1}\) & (3) ARMAX(1,1)\(+r_{t-1}\) \\
\midrule\noalign{}
\endhead
\bottomrule\noalign{}
\endlastfoot
Constant & 0.0721 & 0.1008 & 0.0580 \\
SE(Constant) & 0.0171 & 0.0170 & 0.0134 \\
AR(1) & -0.2255 & & 0.5384 \\
SE(AR(1)) & 0.1257 & & 0.1923 \\
MA(1) & & & -0.5386 \\
SE(MA(1)) & & & 0.2139 \\
\(r_{t-1}\) & & -0.3834 & -0.3903 \\
SE(\(r_{t-1}\)) & & 0.0902 & 0.1071 \\
\(R^2\) & 0.0509 & 0.2315 & 0.2811 \\
Adj. \(R^2\) & 0.0351 & 0.2186 & 0.2439 \\
\end{longtable}
}

\subsubsection{Paper Table 4}\label{paper-table-4}

{\def\LTcaptype{none} % do not increment counter
\begin{longtable}[]{@{}lrrr@{}}
\toprule\noalign{}
& (1) AR(1) & (2) \(r_{t-1}\) & (3) ARMAX(1,1)\(+r_{t-1}\) \\
\midrule\noalign{}
\endhead
\bottomrule\noalign{}
\endlastfoot
Constant & 0.0605 & 0.1007 & 0.0594 \\
SE(Constant) & 0.0127 & 0.0170 & 0.0146 \\
AR(1) & -0.2214 & & 0.5594 \\
SE(AR(1)) & 0.1255 & & 0.1517 \\
MA(1) & & & -0.5688 \\
SE(MA(1)) & & & 0.2080 \\
\(r_{t-1}\) & & -0.3717 & -0.3784 \\
SE(\(r_{t-1}\)) & & 0.0904 & 0.1250 \\
\(R^2\) & 0.0501 & 0.2227 & 0.2779 \\
Adj. \(R^2\) & 0.0340 & 0.2095 & 0.2406 \\
\end{longtable}
}

\subsubsection{Side-by-side comparison}\label{side-by-side-comparison}

{\def\LTcaptype{none} % do not increment counter
\begin{longtable}[]{@{}
  >{\raggedright\arraybackslash}p{(\linewidth - 12\tabcolsep) * \real{0.1111}}
  >{\raggedleft\arraybackslash}p{(\linewidth - 12\tabcolsep) * \real{0.1481}}
  >{\raggedleft\arraybackslash}p{(\linewidth - 12\tabcolsep) * \real{0.1481}}
  >{\raggedleft\arraybackslash}p{(\linewidth - 12\tabcolsep) * \real{0.1481}}
  >{\raggedleft\arraybackslash}p{(\linewidth - 12\tabcolsep) * \real{0.1481}}
  >{\raggedleft\arraybackslash}p{(\linewidth - 12\tabcolsep) * \real{0.1481}}
  >{\raggedleft\arraybackslash}p{(\linewidth - 12\tabcolsep) * \real{0.1481}}@{}}
\toprule\noalign{}
\begin{minipage}[b]{\linewidth}\raggedright
\end{minipage} & \begin{minipage}[b]{\linewidth}\raggedleft
Rep.~(1)
\end{minipage} & \begin{minipage}[b]{\linewidth}\raggedleft
Paper (1)
\end{minipage} & \begin{minipage}[b]{\linewidth}\raggedleft
Rep.~(2)
\end{minipage} & \begin{minipage}[b]{\linewidth}\raggedleft
Paper (2)
\end{minipage} & \begin{minipage}[b]{\linewidth}\raggedleft
Rep.~(3)
\end{minipage} & \begin{minipage}[b]{\linewidth}\raggedleft
Paper (3)
\end{minipage} \\
\midrule\noalign{}
\endhead
\bottomrule\noalign{}
\endlastfoot
Constant & 0.0721 & 0.0605 & 0.1008 & 0.1007 & 0.0580 & 0.0594 \\
SE(Constant) & 0.0171 & 0.0127 & 0.0170 & 0.0170 & 0.0134 & 0.0146 \\
AR(1) & -0.2255 & -0.2214 & & & 0.5384 & 0.5594 \\
SE(AR(1)) & 0.1257 & 0.1255 & & & 0.1923 & 0.1517 \\
MA(1) & & & & & -0.5386 & -0.5688 \\
SE(MA(1)) & & & & & 0.2139 & 0.2080 \\
\(r_{t-1}\) & & & -0.3834 & -0.3717 & -0.3903 & -0.3784 \\
SE(\(r_{t-1}\)) & & & 0.0902 & 0.0904 & 0.1071 & 0.1250 \\
\(R^2\) & 0.0509 & 0.0501 & 0.2315 & 0.2227 & 0.2811 & 0.2779 \\
Adj. \(R^2\) & 0.0351 & 0.0340 & 0.2186 & 0.2095 & 0.2439 & 0.2406 \\
\end{longtable}
}

\subsubsection{Replicated Table 4 (1946 -
2024)}\label{replicated-table-4-1946---2024}

{\def\LTcaptype{none} % do not increment counter
\begin{longtable}[]{@{}lrrr@{}}
\toprule\noalign{}
& (1) AR(1) & (2) \(r_{t-1}\) & (3) ARMAX(1,1)\(+r_{t-1}\) \\
\midrule\noalign{}
\endhead
\bottomrule\noalign{}
\endlastfoot
Constant & 0.0749 & 0.1009 & 0.0584 \\
SE(Constant) & 0.0154 & 0.0147 & 0.0116 \\
AR(1) & -0.2739 & & 0.3973 \\
SE(AR(1)) & 0.1096 & & 0.1920 \\
MA(1) & & & -0.4241 \\
SE(MA(1)) & & & 0.2285 \\
\(r_{t-1}\) & & -0.4081 & -0.4163 \\
SE(\(r_{t-1}\)) & & 0.0765 & 0.1005 \\
\(R^2\) & 0.0750 & 0.2698 & 0.3006 \\
Adj. \(R^2\) & 0.0630 & 0.2604 & 0.2726 \\
\end{longtable}
}

    \subsection{Conclusion}\label{conclusion}

The overall replication was successful. The small differences were
likely due to the paper using data from CRSP while us using data from
WRDS.

    

    

    

    

    


    % Add a bibliography block to the postdoc
    
    
    
\end{document}
